\section{Experiments and Results}\label{sec:experiments}

% One subsection per experiment: question, setup, table, figure, what it shows.
% Tables and figures are read from the current result set in results/<experiment>/ by \experimenttable and \experimentfigure;
% the first argument is the config name in experiments/ (without .yaml). See report/README.md.

Each experiment varies one sampling choice, sometimes together with the number of particles, and keeps everything
else fixed. The tables give the mean and standard deviation over 20 seeds; the figures show the mean with its 95\%
confidence interval. A difference is called clear when the 95\% intervals do not overlap (\cref{sec:method}).

\subsection{Resampling scheme}\label{sec:e01}

\paragraph{Question and setup.} Do the four resampling schemes give different accuracy or particle diversity when
everything else is equal? The localization filter uses the motion-model proposal, resamples at every step so that only
the scheme differs, and runs with 50, 200 and 1000 particles. The same comparison in FastSLAM 1.0 uses 10 and 50
particles. The expectation was that the lower-variance schemes keep more distinct particles and give slightly lower
error than multinomial resampling, most visibly with few particles.

\paragraph{Results.} In localization (\cref{tab:E01_resampling_scheme_localization}) the schemes are equally accurate.
At every particle count their position errors differ by less than the seed-to-seed spread; with 1000 particles all four
reach 0.13\,m, and at 200 particles 0.25 to 0.27\,m. They differ clearly in diversity. After a resampling step,
systematic resampling keeps 72\% distinct particles with 1000 particles, stratified 67\%, residual 65\% and multinomial
54\%, and the same order holds at 50 and 200 particles. The effective sample size before resampling does not depend on
the scheme, since it is computed before the scheme acts.

In FastSLAM 1.0 (\cref{tab:E01_resampling_scheme_slam}) the order of diversity is the same, from 86\% distinct particles
for systematic to 61--62\% for multinomial resampling, and here it carries over to accuracy. With 10 particles
systematic resampling gives a clearly lower trajectory error than multinomial resampling (ATE 0.23 against 0.32\,m) and
a clearly better map (0.105 against 0.139\,m). With 50 particles stratified resampling gives the best map, clearly
better than multinomial (0.069 against 0.097\,m).

The runtimes of the four schemes reflect the reused implementations rather than the schemes themselves. The
multinomial resampler searches the cumulative weights linearly for every new particle, and the stratified resampler
recomputes the sum of all weights for every particle, so both take $O(N^2)$ time; the residual resampler passes its
remaining particles to the same multinomial search. The systematic resampler makes one pass over the weights, $O(N)$.
All four schemes can be implemented in $O(N)$, so only the accuracy and diversity results compare the schemes; the
runtime column compares these implementations.

\experimenttable{E01_resampling_scheme_localization}{Resampling schemes in localization, resampling at every step.}
\experimentfigure{E01_resampling_scheme_localization}{metrics}{Resampling schemes in localization: metrics against the number of particles.}
\experimenttable{E01_resampling_scheme_slam}{Resampling schemes in FastSLAM 1.0, resampling at every step.}
\experimentfigure{E01_resampling_scheme_slam}{metrics}{Resampling schemes in FastSLAM 1.0: metrics against the number of particles.}

\subsection{When to resample}\label{sec:e02}

\paragraph{Question and setup.} Should the filter resample at every step, never, or only when the weights have
degenerated, and how much does the threshold matter? The localization filter uses 100 particles, the motion-model
proposal and systematic resampling. The adaptive rules resample when $\mathrm{ESS} < \tau N$ with $\tau$ = 0.2, 0.5 or
0.8, or when $1 / \max_i \w^{(i)} < \tau N$ with $\tau$ = 0.1, 0.2 or 0.5. The expectation was that never resampling
loses track, resampling at every step wastes diversity, and an ESS threshold around 0.5 does best.

\paragraph{Results.} Never resampling fails (\cref{tab:E02_when_to_resample}): the effective sample size falls to about
1\% of the particles and the position error, 1.39\,m, is three times that of any rule that resamples. The median NEES
of $1.6 \times 10^9$ shows that its estimate is confidently wrong. All seven rules that resample reach the same accuracy,
0.45 to 0.49\,m, with differences far inside the seed-to-seed spread, and the same median NEES of 43 to 51. They
differ in how often they resample: the thresholds of 0.2 resample at about a third of the steps (34\% for the ESS rule,
35\% for the largest-weight rule), the middle thresholds at about half (54\% and 48\%), and the highest at 75--80\%.
Resampling less often leaves fewer distinct particles after each resampling step, 25\% at the lowest thresholds against
64\% when resampling at every step, because the weights are more degenerate when resampling happens. With this proposal
and world, the threshold therefore changes the cost and the diversity per resampling step, but not the accuracy; the
ESS rule and the largest-weight rule behave almost identically at comparable thresholds.

\experimenttable{E02_when_to_resample}{Rules for when to resample, localization with 100 particles.}
\experimentfigure{E02_when_to_resample}{metrics}{When to resample: metrics per rule.}
\experimentfigure{E02_when_to_resample}{trace_ess}{When to resample: effective sample size over time.}

\subsection{Proposal distribution}\label{sec:e03}

\paragraph{Question and setup.} Do proposals that use the current measurement beat the motion model, and at which
particle counts? In localization the motion model, the auxiliary particle filter and the extended Kalman particle filter
run with 50, 100, 300 and 1000 particles, systematic resampling and an ESS threshold of 0.5. In SLAM, FastSLAM 1.0 and
2.0 run with 5, 10, 25 and 100 particles, systematic resampling and the upstream rule, resampling when
$\mathrm{ESS} < N / 1.5$. The expectation was that measurement-informed proposals reach a given error with fewer
particles at a higher cost per particle.

\paragraph{Localization.} The auxiliary particle filter keeps the largest effective sample size at every particle count,
67--83\% against 44--55\% for the motion model (\cref{tab:E03_proposal_localization}), and from 100 particles on it is
the most accurate: 0.26\,m against 0.46\,m for the motion model at 100 particles, and 0.12 against 0.13\,m at 1000.
Its runtime per step is about twice that of the motion model with few particles and similar with 1000 particles. The
extended Kalman particle filter keeps only about 20\% effective particles at every particle count. Its error falls with
more particles, from 0.59\,m at 50 to 0.25\,m at 1000, and with 50 particles it is as accurate as the other two, but
with 1000 particles it is clearly worse than both (0.25 against 0.12 and 0.13\,m) and six times slower. In this
setting its linearized proposal does not make the weights more even, which is the point of a measurement-informed
proposal (\cref{sec:discussion}).

\paragraph{SLAM.} FastSLAM 2.0 keeps a clearly larger effective sample size than FastSLAM 1.0 at every particle count,
79--81\% against 75--76\% (\cref{tab:E03_proposal_slam}). Its advantage in accuracy is largest with few particles. With
5 particles it gives a clearly lower trajectory error (ATE 0.26 against 0.34\,m) and a clearly better map (0.10 against
0.16\,m); with 10 particles the trajectory error is still clearly lower (0.22 against 0.26\,m). With 25 particles
FastSLAM 1.0 has the clearly lower trajectory error (0.21 against 0.28\,m), and with 100 particles the two are within the
spread. FastSLAM 2.0 costs 43--55\% more time per step. With 25 particles FastSLAM 1.0 builds a map as good as FastSLAM
2.0 with 5 to 10 particles, at a similar runtime.

\experimenttable{E03_proposal_localization}{Proposal distributions in localization.}
\experimentfigure{E03_proposal_localization}{metrics}{Proposal distributions in localization: metrics against the number of particles.}
\experimenttable{E03_proposal_slam}{FastSLAM 1.0 and 2.0.}
\experimentfigure{E03_proposal_slam}{metrics}{FastSLAM 1.0 and 2.0: metrics against the number of particles.}

\subsection{Number of particles}\label{sec:e04}

\paragraph{Question and setup.} How do error, particle diversity and runtime change with the number of particles, and
where do more particles stop helping? The localization filter uses the motion-model proposal, systematic resampling and
an ESS threshold of 0.5, with 10 to 2500 particles.

\paragraph{Results.} The error falls steeply up to a few hundred particles and then levels off
(\cref{tab:E04_particle_count}): from 1.50\,m with 10 particles to 0.46\,m with 100, 0.21\,m with 250 and 0.16\,m with 500,
and then only to 0.13\,m with 1000 and 0.12\,m with 2500. The spread over seeds shrinks in the same way, from 1.1\,m
to 0.02\,m: with few particles some seeds do not find the true pose within the 30 steps after the uniform start, and
with many particles all seeds do. The median NEES tells the same story. It is above $10^6$ with 10 and 25 particles, where the particle cloud
has collapsed onto a wrong pose, and falls below 3 from 250 particles on. Values below 3 mean the particle spread is
somewhat larger than the actual error, which fits the filter's assumed noise being larger than the true noise. The
runtime per step grows from 0.18\,ms to 3.7\,ms. For this world, about 250 to 500 particles give most of the achievable
accuracy at a small fraction of the cost of 2500.

\experimenttable{E04_particle_count}{Number of particles, localization with the motion-model proposal.}
\experimentfigure{E04_particle_count}{metrics}{Accuracy and runtime against the number of particles.}

\subsection{Learning the noise parameters}\label{sec:e11}

\paragraph{Question and setup.} The experiments above take the filter's noise settings as given. Can they be
estimated from a recorded run instead? Each of the four components of $\theta$ is swept in turn while the other
three are held at the values the world actually uses, and every value is scored by the log likelihood
\eqref{eq:evidence} of the recorded measurements; the highest score gives the maximum-likelihood estimate. The
filter uses the motion-model proposal, systematic resampling, an ESS threshold of 0.5 and 2000 particles, and
starts around the true initial pose. The world is unchanged, so the answers are known: 0.2\,m and 0.05\,rad for the
measurements, 0.005\,m and 0.002\,rad for the motion.

\paragraph{The measurement parameters are recovered.} The sweeps of the assumed range and bearing noise have a
clear maximum at the true value (\cref{tab:E11_learn_measurement_range,tab:E11_learn_measurement_bearing}). The
best range noise is 0.19\,m against a true 0.2\,m, and the best bearing noise 0.047\,rad against a true
0.05\,rad, both within one step of the sweep. The shape on either side is what \cref{sec:learning} predicts.
Assuming too little noise makes the measurements that actually arrived look impossible and the log likelihood
falls steeply, from 205 at 0.19\,m to $-252$ at 0.08\,m. Assuming too much noise makes every pose explain the
measurements equally well and the score decays gently, to 98 at 0.8\,m. The penalty is strongly one-sided:
halving the assumed range noise costs about 450 in log likelihood, while quadrupling it costs about 100, the
asymmetry that \eqref{eq:gaussian-score} predicts. A filter whose noise is set by hand is therefore much safer
overstating it than understating it, which is what the upstream settings do.

On this world the most plausible setting is also close to the most accurate one. The position RMSE is lowest,
0.063\,m, at 0.19 and 0.253\,m of assumed range noise, the same place as the likelihood maximum, so learning the
model and tuning for accuracy do not conflict here.

\paragraph{The sweep matches the closed form.} The measured curve can be checked against
\eqref{eq:gaussian-score}, which gives the expected score of one residual, relative to its maximum, as
$-\log(\sigma/\sigma^*) - (\sigma^{*2}/2\sigma^{2} - 1/2)$. The range sweep has four landmarks over 30 steps,
so 120 such residuals, and nothing is fitted. Relative to the peak, the prediction at 0.142, 0.253, 0.45 and
0.8\,m is $-18$, $-6$, $-49$ and $-110$, against measured $-20$, $-5$, $-47$ and $-108$. The one place it parts
company is the far left: at 0.08\,m the prediction is $-205$ and the measurement $-457$, because there the
weights degenerate and the downward bias of the estimator, which \eqref{eq:gaussian-score} does not describe,
dominates. The agreement is evidence that the sweep measures what it is meant to, and the disagreement marks
where the particle approximation stops being trustworthy.

\paragraph{The motion parameters are not identifiable from 30 steps.} The same sweeps over the assumed forward
and turn noise are flat (\cref{tab:E11_learn_motion_forward,tab:E11_learn_motion_turn}): the log likelihood stays
at 205 to 206 while the assumed forward noise is varied by a factor of 30 and the assumed turn noise by a factor
of 30. The true values, 0.005\,m and 0.002\,rad over 30 steps, are small next to the measurement noise that the
filter corrects with at every step, so the recorded measurements barely depend on them. The parameters are in the
model, and the data say almost nothing about them.

This is a property of the data and not of the method. Repeating the turn sweep on a trajectory of 400 steps, with
nothing else changed, turns the flat curve into one with a maximum at the true value
(\cref{fig:E11_learn_motion_turn_long-metrics}): 0.0025\,rad against a true 0.002\,rad, and an assumed noise of
0.0008\,rad now costs about 8000 in log likelihood. Because the parameters are tied across time slices, every
extra step is evidence about the same number, and by \eqref{eq:fisher} the curvature of the likelihood grows with
the number of steps, so enough steps resolve what 30 steps cannot.
The one-sidedness is even stronger here: too small an assumed drift is rejected by thousands of nats, while ten
times too much costs about 30.

\paragraph{Independent recordings agree exactly when the parameter is identifiable.} Each seed is a separate
recording of the robot, so the estimate can be formed from each seed on its own and the results compared. Where
the sweep has a maximum the recordings agree: all 20 recordings pick 0.047\,rad for the bearing noise, and 17 of
20 pick 0.19\,m for the range noise with the other three picking the next value up. Where the sweep is flat they
do not: the 20 estimates of the turn noise from 30-step runs are spread over the whole sweep, from 0.0008 to
0.025\,rad, with no value chosen by more than a few recordings. On 400 steps they concentrate again, 9 of 20 on
0.0025\,rad and all within 0.0008 to 0.0045\,rad. Agreement between recordings is thus a usable test of whether a
parameter is being estimated at all, and it also tests the assumption behind the tying: a noise level that changed
between runs would show as a wide scatter even where the likelihood is sharply peaked.

\paragraph{Learning needs the filter to track.} Repeating the range sweep from the uniform start used in
\crefrange{sec:e01}{sec:e04} removes the maximum entirely (\cref{fig:E11_start_regime-metrics}). With the true, small
noise the filter cannot find the robot within 30 steps from a uniform prior: the position RMSE is 2.3\,m against
0.063\,m from a known start, and the log likelihood is $-1.3 \times 10^{4}$ against 205. The score then keeps
improving as the assumed noise is widened, to $-5.2 \times 10^{3}$ at 0.8\,m, because a wider noise lets the
particles spread enough to be useful, and the estimates from individual recordings scatter across the whole sweep,
with only 2 of 20 agreeing. What the likelihood reports in this regime is which setting lets this filter work, not
which setting generated the data.

\paragraph{How much inference learning needs.} Repeating the range sweep with 25 to 2000 particles shows the
transition (\cref{fig:E11_estimator_particles-best_per_seed}). At the true value the estimated log likelihood rises from
$28 \pm 316$ with 25 particles to $184 \pm 65$ with 100, $203 \pm 16$ with 500 and $205 \pm 13$ with 2000: the
downward bias and the spread of the estimator both shrink as the particle approximation improves, as
\cref{sec:learning} anticipates. The estimate follows. With 25 particles no two recordings agree and the median
estimate is 0.22\,m; with 100 particles 12 of 20 recordings pick 0.19\,m, with 500 particles 14 of 20 and with
2000 particles 17 of 20. The resampling scheme, in contrast, makes almost no difference
(\cref{tab:best-E11_estimator_resampler}): with 200 particles all four schemes put the maximum at 0.19\,m and 14 or 15
of 20 recordings agree, matching \cref{sec:e01}, where the schemes differed in particle diversity rather than in
accuracy.

\experimentfigure{E11_learn_measurement_range}{metrics}{Sweeping the assumed range noise: the log likelihood peaks at the value the world uses, 0.2\,m, and the position error is lowest at the same place.}
\experimenttable{E11_learn_measurement_range}{Sweeping the assumed range noise; the world uses 0.2\,m.}
\experimenttable{E11_learn_measurement_bearing}{Sweeping the assumed bearing noise; the world uses 0.05\,rad.}
\experimenttable{E11_learn_motion_forward}{Sweeping the assumed forward motion noise; the world uses 0.005\,m.}
\experimenttable{E11_learn_motion_turn}{Sweeping the assumed turn noise over 30 steps; the world uses 0.002\,rad.}
\experimentfigure{E11_learn_motion_turn_long}{metrics}{The same turn sweep over 400 steps: enough steps identify a value that 30 steps cannot.}
\experimentfigure{E11_start_regime}{metrics}{The range sweep from a known start and from the uniform start of \crefrange{sec:e01}{sec:e04}. Note the scale: from the uniform start the estimate is four orders of magnitude lower and has no maximum.}
\experimentfigure{E11_estimator_particles}{best_per_seed}{The range noise that each recording picks out on its own, against the number of particles.}
\experimentbest{E11_estimator_resampler}{The range noise picked out under each resampling scheme, with 200 particles.}
